\documentclass[12pt]{article}
\usepackage[utf8]{inputenc}
\usepackage[T1]{fontenc}
\usepackage{hyperref}
\title{Kalanadi Language Resources for K-12 Curriculum Development}
\author{Compiled by Claude Sonnet 5.5 (claude-sonnet-5-5) for Andr\'as Kornai}
\date{}
\begin{document}
\maketitle

\section{Introduction}
Kalanadi (ISO 639-3 \texttt{wkl}, Glottocode kala1255) is a Southern Dravidian language of India. Wikipedia cites about 750 native speakers (2004 survey) and 88\% lexical similarity with Pathiya. The search found very little about it.

\section{Primary Educational Resources}
No primers, textbooks or school programmes were found.

\section{Linguistic and Reference Materials}
\begin{itemize}
\item Wikipedia, Kalanadi language: \url{https://en.wikipedia.org/wiki/Kalanadi_language}
\item Glottolog languoid page for Kalanadi: \url{https://glottolog.org/resource/languoid/id/kala1255.bigmap.html}
\item Wikipedia, Pathiya language (closest relative): \url{https://en.wikipedia.org/wiki/Pathiya_language}
\end{itemize}
No dictionary or grammar was found.

\section{Religious and Community Resources}
None found.

\section{Digital and Audio Resources}
None found.

\section{Curriculum Development Context}
The language is not known to be taught in schools. Any curriculum would have to start with basic documentation, and could draw on the regional majority language and on related Dravidian languages for comparison.
\end{document}